\begin{lemma}[One-block canonical peripheral spectrum]
    \label{lem:pgvwc07_unital_canonical_spectrum}
    \lean{MPSTensor.isIrreducible_transferMap_of_unital_canonical,
        MPSTensor.exists_peripheral_generator_of_unital_canonical,
        MPSTensor.peripheral_genEigenspace_finrank_eq_one_of_unital_canonical}
    \leanok
    Let $D>0$ and let $A=(A^i)$ satisfy the one-block canonical conditions
    of \cite[Theorem~4]{PerezGarcia2007Matrix}:
    \begin{align}
        \sum_i A^i(A^i)^\dagger        & =\Id,     &
        \sum_i(A^i)^\dagger\Lambda A^i & =\Lambda, & \Lambda & >0,
        \label{eq:pgvwc07_unital_canonical_conditions}
    \end{align}
    and every fixed matrix of $\E_A$ is a scalar multiple of $\Id$.
    Diagonality of $\Lambda$ is unnecessary for the conclusions below.
    Then $\E_A$ is irreducible. If its peripheral spectrum has cardinality
    $p$, then $p>0$ and it is precisely
    $\{\gamma^k:0\leq k<p\}$ for a primitive $p$-th root $\gamma$.
    Each peripheral maximal generalized eigenspace has dimension one.
    Thus counting distinct peripheral eigenvalues also counts their
    algebraic multiplicities.
\end{lemma}

\begin{proof}
    \leanok
    Unitality fixes $\Id$ and gives spectral radius one. Together with the
    faithful adjoint fixed point $\Lambda$ and uniqueness of positive fixed
    matrices up to scalars, these are the spectral characterization of
    irreducibility. The irreducible unital peripheral spectrum is cyclic.
    Its primitive-root powers are distinct, so its cardinality identifies
    the cyclic order with $p$. Positivity and unitality exclude nontrivial
    peripheral Jordan blocks; the irreducible peripheral eigenspaces are
    one-dimensional. These are the peripheral spectral facts used in
    \cite[proof of Theorem~5]{PerezGarcia2007Matrix}.
\end{proof}

\begin{theorem}[Periodic decomposition from the unital canonical form]
    \label{thm:pgvwc07_unital_periodic_decomposition}
    \lean{MPSTensor.pgvwc07_unital_stateVector_decomposition}
    \leanok
    \uses{def:pgvwc07_projector_component_chain,
        lem:pgvwc07_unital_canonical_spectrum}
    Under the hypotheses of
    Lemma~\ref{lem:pgvwc07_unital_canonical_spectrum}, let $p$ count the
    peripheral eigenvalues. There are nonzero orthogonal projections
    $(Q_u)_{u\in\mathbb Z/p\mathbb Z}$ satisfying
    \begin{align}
        \sum_uQ_u & =\Id, & Q_uA^i & =A^iQ_{u+1}.
        \label{eq:pgvwc07_unital_projector_shift}
    \end{align}
    For every positive ring length $N$, define the component chains by
    $B_{u,j}^i=Q_{u+j}A^iQ_{u+j+1}$. Each chain repeats a $p$-site pattern
    and has bond dimension $D$, and
    \begin{align}
        \ket{V^{(N)}(A)} & =\sum_{u\in\mathbb Z/p\mathbb Z}
        \ket{V^{(N)}(B_u)}.
        \label{eq:pgvwc07_unital_periodic_decomposition}
    \end{align}
    If $p\nmid N$, every component vector is zero, hence so is the original
    vector. This is \cite[Theorem~5]{PerezGarcia2007Matrix}, directly in
    its unital orientation.
\end{theorem}

\begin{proof}
    \leanok
    \uses{lem:pgvwc07_unital_canonical_spectrum, lem:kraus_cyclic_shift,
        thm:pgvwc07_boundary_trace_coefficients}
    The cyclic peripheral decomposition supplies orthogonal projectors with
    $\E_A(Q_{u+1})=Q_u$. Their sum is the identity, and cyclicity makes each
    nonzero. The Kraus shift gives
    (\ref{eq:pgvwc07_unital_projector_shift}) for the original matrices
    $A^i$. Inserting the resolution of the identity into the trace proves
    (\ref{eq:pgvwc07_unital_periodic_decomposition}). At a non-divisible
    length, the starting and final projectors differ; their orthogonality
    and cyclicity of the trace kill each component.
\end{proof}
